\documentclass[11pt]{article}
\usepackage[margin=1in]{geometry}
\usepackage{amsmath,amssymb,amsthm}
\usepackage{booktabs}
\usepackage{hyperref}
\usepackage[numbers,sort&compress]{natbib}
\hypersetup{hidelinks}
\emergencystretch=3em

\newtheorem{theorem}{Theorem}
\newtheorem{proposition}{Proposition}
\newtheorem{corollary}{Corollary}
\newtheorem{definition}{Definition}
\newtheorem{remark}{Remark}
\newtheorem{conjecture}{Conjecture}
\newtheorem{observation}{Observation}

\newcommand{\Hyp}{\Xi}
\newcommand{\feas}{\mathrm{feas}}
\newcommand{\Sec}{\mathcal S}
\newcommand{\rectrho}{\rho}
\newcommand{\obs}{\mathrm{obs}}
\newcommand{\Hech}{\check H}

\title{The Cohomological Layer of Set-Theoretic Estimation:\\
What Gluing Buys, What It Cannot, and Where the Speedups Live\\
\large An exploratory research memo on the \v{C}ech / sheaf machinery of
project \texttt{thejeb}}
\author{Project \texttt{thejeb}}
\date{July 28, 2026}

\begin{document}
\maketitle

\begin{abstract}
This is an \emph{exploratory} memo, not a results paper: it mixes a
survey of machine-checked Lean~4 theorems with clearly marked
conjectures, proof sketches, and open directions, and the two kinds of
statement are kept typographically and verbally separate throughout.
The subject is the sheaf-cohomological layer of the project's
mechanization of set-theoretic estimation (STE)
\citep{combettes1993}: the constraint-side sheaf, the variable-side
section presheaf and its \v{C}ech gluing obstruction
\citep{abramsky2011sheaf,abramsky2012cohomology}, the constant- and
twisted-coefficient cochain complexes, the quantitative obstruction
counts, and the rectangle-cover complexity $\rectrho$.  We ask two
questions.  First, can the cohomological view \emph{speed up} STE ---
help avoid the $2^N$ feasible-set blow-up, decompose, glue, or certify?
Second, how does it \emph{connect} to the rest of the development ---
treewidth, bucket elimination, junction trees, adaptive consistency,
the dynamic-frame corpus gluing, and Carlson's cryptanalytic
instantiation \citep{carlson2012}?  The centerpiece of the exploration
is the \emph{support--cover gluing theorem}, conjectured during this
exploration and since \emph{mechanized} (sorry-free, module
\path{Ste.SupportCover}): over any cover of the variable set, the
\v{C}ech gluing condition for the genuine section presheaf holds
\emph{iff} the feasible set is generated by constraints whose scopes
refine the cover.  It subsumes two previously mechanized special
cases, collapses several open questions in the repository's
docstrings, and, we argue, relocates the entire computational
difficulty of STE into \emph{computing local sections}, with gluing
itself free; its downstream corollaries (the scope lower bound, the
gluing-width invariant, the arity-vs-treewidth separation) remain
clearly marked sketches.  From that vantage point we
identify three concrete speedup routes (one essentially proven, two
conjectural), explain why the na\"ive hope ``a cohomological invariant
of the cover detects hardness'' is refuted by the machine-checked
acyclicity of the constant-coefficient complex, and rank the
mechanization targets that would buy the most.
\end{abstract}

\section{Orientation: three localities, one wall}
\label{sec:orientation}

Everything in this memo is about one wall.  An STE problem
\citep{combettes1993} is a family of property sets $S_i \subseteq
\Hyp$ and its answer object is the feasibility set $T = \bigcap_i S_i$
(\path{feasibilitySet}, \path{mem_feasibilitySet} in
\path{Ste.Basic}).  When $\Hyp = \prod_{v \in V} A_v$ is a product
over $N$ variables, $T$ lives in a space of size exponential in $N$,
and the project has machine-checked that the wall is real: crisp
property sets can realize \emph{any} of the $2^N$ extensional states
(\path{card_allExactStates},
\path{arbitrary_subset_as_single_property} in
\path{Ste.DynamicFrame.FiniteSolver}).  The companion note on
feasible-set shrinking catalogues the answer-preserving reductions;
this memo examines the specifically \emph{cohomological} machinery and
asks what it contributes to getting around the wall.

The development contains three distinct axes of locality, and it is
worth fixing them at the outset because the cohomology lives on
exactly one of them.

\begin{enumerate}
\item \textbf{The constraint axis} (index set $I$).  The assignment $J
  \mapsto \feas(J) = \bigcap_{i \in J} S_i$ from constraint subsets to
  partial feasible sets sends unions to intersections
  (\path{partialFeasibilitySet_iUnion}), so any cover of the
  constraint index recovers the global feasible set as the limit of
  the local ones (\path{feasibilitySet_eq_iInter_cover}, in
  \path{Ste.Sheaf}).  \emph{This side is a sheaf for free}: colimits
  of constraints go to limits of feasible sets, and no obstruction can
  live here.  (Under the formal-concept reading of
  \path{Ste.HyperFrame}, $\feas$ is a lower polar of the satisfaction
  context \citep{ganter1999fca}
  (\path{partialFeasibilitySet_eq_lowerPolar}), and polars always
  turn joins into meets.)
\item \textbf{The variable axis} (contexts $W \subseteq V$).  For a
  constraint $T \subseteq \prod_v A_v$ and a context $W$, the
  \emph{local sections} $\Sec_T(W) = $ \path{localSections}$\,T\,W$
  are the partial assignments on $W$ that extend to a global solution
  of $T$ --- the image of $T$ under coordinate restriction
  (\path{Ste.VariablePresheaf}).  Restriction along $W' \subseteq W$
  makes $W \mapsto \Sec_T(W)$ a genuine contravariant functor
  (\path{restrict_restrict}, \path{localSections_restrict₂}).  This
  presheaf is \emph{not} in general a sheaf, and the failure is
  precisely variable coupling: this is where all the \v{C}ech
  machinery of the repository lives, and where this memo spends most
  of its time.
\item \textbf{The corpus axis} (document sets $D$, in
  \path{Ste.DynamicFrame}).  The feasible-world construction $D
  \mapsto \mathrm{feasible}(D)$ is contravariant in the active corpus
  with an \emph{exact} binary gluing law: a hypothesis is feasible on
  $D \cup E$ iff it is feasible on each piece and the
  \emph{cross-constraint obstruction} --- the set of constraints
  supported only in the union and violated by the hypothesis ---
  is empty (\path{crossObstruction}, \path{mem_feasible_union_iff}
  in \path{Ste.DynamicFrame.PresheafGluing}); a global extension of a
  compatible pair of local worlds exists iff that obstruction
  vanishes, and is then unique
  (\path{globalExtension_nonempty_iff},
  \path{globalExtension_subsingleton}, \path{glue}).
\end{enumerate}

The slogan the mechanization supports: \emph{the constraint side glues
for free, the corpus side glues with an exact and computable
obstruction term, and the variable side is where genuine
cohomology-shaped failure occurs}.  Sections~\ref{sec:inventory}
and~\ref{sec:coefficients} survey what is machine-checked on the
variable side.  Section~\ref{sec:supportcover} states the
support--cover gluing theorem --- conjectured during this exploration
and since mechanized in \path{Ste.SupportCover} --- that we believe
organizes all of it.  Sections~\ref{sec:algorithms}--\ref{sec:quant} draw the
algorithmic and quantitative consequences, and
Section~\ref{sec:mechanization} ranks what to mechanize next.

\section{Inventory: the machine-checked variable-side story}
\label{sec:inventory}

All statements in this section are proved in Lean, sorry-free, in the
named files; we cite theorem names exactly.

\subsection{The singleton cover: gluing iff rectangular}

Call $T$ \emph{rectangular} if $T = \prod_v P_v$ for per-variable
sides $P_v$ (\path{IsRectangle}).  Rectangular feasibility problems
are the tractable fragment: the feasible set of a rectangular family
is itself a rectangle (\path{rectangular_feasibilitySet}) of
cardinality $\prod_v |{\bigcap_i P_{i,v}}|$
(\path{rectangular_encard}), representable in $\sum_v$ space.  The
two-bit coupling $\Delta = \{f \mid f(0) = f(1)\}$
(\path{diagonal}) is provably not rectangular
(\path{diagonal_not_rectangular}), has exactly two solutions
(\path{diagonal_encard}), and every rectangle containing it is the
whole four-point space
(\path{rectangle_superset_diagonal_encard}): the tightest
per-variable view of a coupling carries zero information.

For the cover of $V$ by singletons, the \v{C}ech data is computed
exactly (\path{Ste.CechObstruction}).  Overlaps of distinct singleton
contexts are empty (\path{singletonCover_overlap_empty}), so a
compatible family is a point of the box $\prod_v \pi_v(T)$
(\path{compatibleFamilies}); a family glues iff, read as an
assignment, it lies in $T$ (\path{glues_iff_mem},
\path{gluedFamilies_eq}).  The obstruction number
\[
  \obs(T) \;=\; |\mathrm{compatible}| - |\mathrm{glued}|
\]
(in Lean, \path{cechObstruction}) vanishes on rectangles
(\path{rectangular_cechObstruction}) and
equals $2$ on the diagonal (\path{diagonal_cechObstruction}), with
the mixed family $(0 \mapsto \mathsf{false},\, 1 \mapsto
\mathsf{true})$ the explicit stuck cocycle
(\path{diagonal_mixed_compatible_not_glues}).  Crucially the
singleton-cover story is closed as an \emph{iff}
(\path{cechVanishes_iff_rectangular}, in
\path{Ste.CouplingLowerBound}):
\[
  \text{every compatible family glues}
  \iff
  T \text{ is a rectangle.}
\]

\subsection{Arbitrary covers}

\path{Ste.CechCover} lifts the gluing level to an arbitrary cover $U
: J \to \mathcal P(V)$ with possibly nonempty overlaps, where the
compatibility condition (agreement of local sections on overlaps) is
substantive (\path{CompatibleFamily}, \path{GluesCover},
\path{CechVanishesCover}).  The headline is a genuine sheaf theorem:
for a rectangular constraint and \emph{any} cover, every compatible
family of local sections glues, uniquely
(\path{rectangular_cechVanishesCover},
\path{rectangular_glueCover_existsUnique}; uniqueness of glues along
any cover is \path{glueCover_unique} and needs no constraint
membership at all).  Vanishing is genuinely cover-dependent: the
trivial one-context cover glues everything
(\path{cechVanishesCover_univ}), the singleton cover is the finest
and hardest case and recovers the closed theory above
(\path{cechVanishesCover_singleton_iff}), and the diagonal witnesses
real failure (\path{diagonal_not_cechVanishesCover}).  The
cover-level count \path{coverCechObstruction} vanishes whenever the
sheaf condition holds (\path{coverCechObstruction_eq_zero}).

\subsection{Scaling the witness: the $n$-fold coupling}

The generalized diagonal $\mathrm{allEqual}(n, \alpha)$ --- all $n$
coordinates equal over an alphabet $\alpha$ (\path{allEqual}) --- has
every margin full (\path{contextSections_allEqual}), so \emph{all}
$|\alpha|^n$ families are compatible
(\path{compatibleFamilies_allEqual_encard}) while only the $|\alpha|$
constant vectors are feasible (\path{allEqual_encard}).  The
obstruction is exactly $|\alpha|^n - |\alpha|$
(\path{cechObstruction_allEqual}), hence at least $2^n - 2$ once
$|\alpha| \ge 2$ (\path{cechObstruction_allEqual_ge}): exponentially
many phantom assignments that any margins-only procedure must
entertain, while the true feasible set keeps constant size.  For $n
\ge 2$ the constraint is genuinely non-rectangular
(\path{allEqual_not_rectangular},
\path{allEqual_not_cechVanishes}).

\section{What the coefficients see}
\label{sec:coefficients}

The Abramsky--Brandenburger program grades gluing failure
cohomologically by linearizing the section presheaf
\citep{abramsky2011sheaf,abramsky2012cohomology}.  The repository
mechanizes two coefficient regimes, and the contrast between them is,
in our view, the single most instructive machine-checked fact in this
layer.

\subsection{Constant coefficients see nothing --- a machine-checked
no-free-lunch}

\path{Ste.CechComplex} builds the constant-coefficient cochain
complex on the full nerve of a cover index --- modules $C^0 = \iota
\to M$, $C^1 = \iota \to \iota \to M$, \dots\ with the standard
alternating coboundaries (\path{cechD0}, \path{cechD1},
\path{cechD2}), the complex identities
(\path{cechComplex_d_comp_d}, \path{cechComplex_d2_comp_d1}), and
honest quotient modules $\Hech^1, \Hech^2$ (\path{cechH1},
\path{cechH2}).  The vanishing theorem
(\path{cechH1_subsingleton}, \path{cechH2_subsingleton}) then proves
$\Hech^1 = \Hech^2 = 0$ over any nonempty index type: the full nerve
is a simplex, and a simplex is acyclic; $\Hech^0$ is exactly the
coefficients (\path{cechH0EquivCoeff}).

The computational moral deserves to be stated bluntly.  \emph{Any
invariant computed from the combinatorics of the cover alone ---
ignoring the section data --- is useless for detecting coupling}: the
machine checks that it is identically zero.  There is no free
topological lunch; whatever a ``cohomological speedup'' means, its
input must include the tables of local sections, i.e.\ exactly the
data the CSP algorithms of
Section~\ref{sec:algorithms} already manipulate.  This is consistent
with the contextuality literature, where the interesting classes live
in relative cohomology with \emph{presheaf} coefficients
\citep{abramsky2012cohomology,caru2017cohomology}.

\subsection{Twisted coefficients see the coupling --- one degree lower
than expected}

\path{Ste.TwistedCech} builds the honest twisted complex: coefficients
in the linearized section presheaf $F_R(W) = (\Sec_T(W)) \to_{\!f} R$
(free $R$-module on genuine local sections, \path{twistedCoeff}),
restriction by pushforward along section restriction
(\path{sectionRes}, \path{twistedRes}, functorial by
\path{twistedRes_twistedRes}), coboundaries \path{twistedD0},
\path{twistedD1} with $d^1 \circ d^0 = 0$
(\path{twisted_d1_comp_d0}), and $\Hech^0$ as \path{twistedH0}.
Two machine-checked results locate the STE coupling precisely.

\begin{enumerate}
\item \emph{The AMB per-section obstruction degenerates in STE}
  (\path{amb_extension_always}).  Abramsky--Mansfield--Barbosa attach
  to a single local section $s_1$ a class $\gamma(s_1) \in \Hech^1$
  vanishing iff $s_1$ extends to a compatible family of the
  linearized presheaf \citep[Prop.~4.2]{abramsky2012cohomology}.  The
  mechanization proves that in STE this extension condition holds for
  \emph{every} section of \emph{every} constraint over \emph{every}
  cover --- because \path{localSections} only ever contains
  restrictions of global solutions, the extension witness is free.
  The $\Hech^1$-graded per-section obstruction is structurally blind
  here: STE's section presheaf bakes global extendability into its
  stalks.  (The long exact sequence and the literal connecting map
  $\gamma$ are \emph{not} mechanized; what is proven is the extension
  condition AMB prove equivalent to $\gamma = 0$.)
\item \emph{The obstruction that fires lives in the cokernel of
  $F_R(V) \to \Hech^0(U, F_R)$}
  (\path{mixedFamily_not_linearlyGlues},
  \path{diagonal_mixed_cocycle_not_in_globalRange}).  A compatible
  family's point-mass cochain is a twisted $0$-cocycle
  (\path{familyCochain_mem_twistedH0}); the family \emph{glues
  $R$-linearly} (\path{LinearlyGlues}) iff that cocycle lifts to a
  formal $R$-combination of global solutions
  (\path{linearlyGlues_iff_mem_range}).  For the two-bit coupling and
  its singleton cover, the mixed family fails to glue linearly over
  \emph{every nontrivial commutative ring} --- negative and rational
  coefficients included.  This is strictly stronger than the
  set-level failure (\path{diagonal_gluing_fails}): in the AMB world
  every no-signalling model has a global section over the signed
  reals \citep{abramsky2011sheaf}, but the STE coupling is not even
  signed-linearly explainable by global data.  Set-level gluing
  always implies linear gluing (\path{GluesCover.linearlyGlues}),
  and rectangles glue linearly over every cover
  (\path{rectangular_linearlyGlues}), so nonmembership is a genuine
  obstruction invariant
  (\path{diagonal_twisted_obstruction_detects}).
\end{enumerate}

Whether the twisted $\Hech^1$ (carried as the submodule pair
\path{twistedCoboundaries1} $\le \ker$ \path{twistedD1}, with
triviality expressed as \path{TwistedH1Trivial}; the literal quotient
type is blocked by a typeclass diamond, documented in the file)
vanishes for pairwise-disjoint covers is stated as \emph{open} in the
source and remains open here.

\section{The support--cover gluing theorem (now mechanized)}
\label{sec:supportcover}

Here is the centerpiece of the exploration.  It occurred to us while
trying to answer the prompt ``does vanishing $\Hech^1$ over a good
cover give an algorithm?''\ and we believe it substantially
reorganizes the layer.  \textbf{Status update: the core proposition
below --- directions (a), (b), and the packaged iff --- has since
been mechanized, sorry-free, in the module \path{Ste.SupportCover}}
(theorem names cited in the statement).  The boundary must stay
crisp, however: \emph{only} Proposition~\ref{prop:supportcover} is
proven; its downstream consequences in this section ---
Corollary~\ref{cor:scopebound}, Definition~\ref{def:gluingwidth}, and
Observation~\ref{obs:arity} --- remain sketches and prose arguments,
not Lean theorems.

Recall \path{HasSupport}$\,T\,\sigma$ (from \path{Ste.Support}):
membership in $T$ depends only on the coordinates in $\sigma$.  Say a
family $(S_i)_{i \in I}$ with supports $\sigma_i$ \emph{generates} $T$
if $T = \bigcap_i S_i$, and say the family's scopes \emph{refine} the
cover $U : J \to \mathcal P(V)$ if every $\sigma_i$ is contained in
some context $U_j$.

\begin{proposition}[support--cover gluing; \textbf{mechanized},
\path{Ste.SupportCover}]
\label{prop:supportcover}
Let $T \subseteq \prod_v A_v$ and let $U : J \to \mathcal P(V)$ be a
cover of the variable set ($\forall v\, \exists j,\ v \in U_j$).
\begin{enumerate}
\item[(a)] If $T$ is generated by a family whose scopes refine $U$,
  then \path{CechVanishesCover}$\,T\,U$: every compatible family of
  genuine local sections glues to a global solution
  (\path{cechVanishesCover_of_refiningGenerators}).
\item[(b)] Conversely, if \path{CechVanishesCover}$\,T\,U$, then $T =
  \bigcap_{j} \pi_{U_j}^{-1}\!\bigl(\Sec_T(U_j)\bigr)$: $T$ is
  generated by the cylinder constraints on its own local sections
  (\path{sectionCylinder}, supported on their contexts by
  \path{hasSupport_sectionCylinder}), a family whose scopes are
  exactly the contexts of $U$; indeed vanishing is \emph{equivalent}
  to this cylinder identity
  (\path{cechVanishesCover_iff_eq_iInter_sectionCylinder}).
\end{enumerate}
Hence the \v{C}ech obstruction of the cover $U$ vanishes \emph{iff}
$T$ admits a generating family whose scopes refine $U$ --- packaged
as the iff ``vanishing $\iff$ there is a per-context generating
family $S_j$ with \path{HasSupport}$\,(S_j)\,(U_j)$''
(\path{cechVanishesCover_iff_exists_scopedGeneration}).
\end{proposition}

\begin{proof}[Proof (exposition; the machine-checked proof is in
\path{Ste.SupportCover})]
(a) Let $(s_j)$ be a compatible family, each $s_j \in \Sec_T(U_j)$
with global witness $f_j \in T$, i.e.\ $f_j|_{U_j} = s_j$, and with
$s_j, s_k$ agreeing on $U_j \cap U_k$.  Choose for each $v$ a context
$j(v) \ni v$ and define $f(v) := s_{j(v)}(v)$.  Compatibility makes
$f(v) = s_j(v)$ for \emph{every} $j$ with $v \in U_j$; in particular
$f|_{U_j} = s_j$ for all $j$.  Now take any generator $S_i$ with
$\sigma_i \subseteq U_{j_0}$.  For $v \in \sigma_i$ we have $f(v) =
s_{j_0}(v) = f_{j_0}(v)$, so $f$ agrees with $f_{j_0}$ on
$\sigma_i$; by \path{HasSupport}, $f \in S_i$ since $f_{j_0} \in S_i$.
As $i$ was arbitrary, $f \in T$, and it restricts to the prescribed
sections.  (Uniqueness of the glue is already mechanized:
\path{glueCover_unique}.)

(b) Each cylinder $\pi_{U_j}^{-1}(\Sec_T(U_j))$ has support $U_j$ and
contains $T$.  Conversely let $f$ lie in the intersection; then $s_j
:= f|_{U_j} \in \Sec_T(U_j)$ for every $j$, and the family $(s_j)$ is
automatically compatible (all restrict the same $f$).  Vanishing
glues it to some $g \in T$ with $g|_{U_j} = f|_{U_j}$ for all $j$;
since $U$ covers $V$, $g = f$ (again \path{glueCover_unique}), so $f
\in T$.
\end{proof}

\begin{remark}[Sanity anchors, all machine-checked]
Direction (b) at $U = $ singletons is exactly
\path{cechVanishes_iff_rectangular}: rectangularity $=$ generation by
singleton-scope constraints.  Direction (a) contains
\path{rectangular_cechVanishesCover} as the special case where the
generators have singleton scopes, which refine \emph{every} cover.
The diagonal's failure at the singleton cover is consistent: its
minimal scope is $\{0,1\}$ (\path{diagonal_support_full}, in
\path{Ste.Support}), contained in no singleton.  And the trivial-cover
success \path{cechVanishesCover_univ} is the case where the single
context contains all scopes.  The proposition also discharges, at the
cover level, the ``residual conjecture'' left explicitly open in the
docstring of \path{Ste.VariablePresheaf} (gluing holds precisely for
the tractable fragments) --- provided ``tractable fragment'' is read
as \emph{scope-refinement}, which per
Observation~\ref{obs:arity} below is emphatically \emph{not} the same
as computational tractability.
\end{remark}

\begin{remark}[Why no acyclicity hypothesis?]
Classical local-to-global results in constraint processing
(tree-clustering, Vorob'ev-style theorems
\citep{dechterpearl1989tree,freuder1982backtrack}) require the cover's
nerve to be acyclic (running intersection).
Proposition~\ref{prop:supportcover} needs nothing of the sort, and
the reason is the same structural fact behind
\path{amb_extension_always}: STE's local sections are \emph{globally
extendable by definition}.  The classical hypotheses are needed when
the local data is merely \emph{locally consistent} (arc- or
$k$-consistent tables) rather than genuinely extendable; see
Section~\ref{sec:algorithms}.  In Abramsky--Brandenburger terms:
empirical models are arbitrary compatible presheaf data and can be
contextual; the presheaf of genuine sections of a global set cannot
be, in the per-section sense --- but family-level gluing still fails
exactly when the cover is too fine for the constraint arity.
\end{remark}

\begin{corollary}[Cohomological scope lower bound; sketch ---
immediate from Proposition~\ref{prop:supportcover} but not itself
mechanized]
\label{cor:scopebound}
If some compatible family over $U$ is stuck (nonvanishing
\path{coverCechObstruction}), then \emph{every} generating family of
$T$ contains a constraint whose scope is contained in no context of
$U$.  In particular the mechanized witnesses
(\path{diagonal_not_cechVanishesCover},
\path{allEqual_not_cechVanishes}) certify, at the singleton cover,
that no unary decomposition of these constraints exists --- and a
single small stuck family over any proposed bag cover certifies that
a bag-local pipeline over those bags is unsound for $T$.
\end{corollary}

\begin{definition}[Gluing width; not mechanized]
\label{def:gluingwidth}
Define $\mathrm{gw}(T)$ as the least $w$ such that
\path{CechVanishesCover}$\,T\,U$ holds for some cover $U$ all of whose
contexts have at most $w+1$ variables.  By
Proposition~\ref{prop:supportcover}, $\mathrm{gw}(T)$ is exactly
(one less than) the minimal \emph{arity} of a generating constraint
family for $T$.
\end{definition}

\begin{observation}[Gluing width is arity, not treewidth]
\label{obs:arity}
$\mathrm{gw}$ must not be confused with treewidth, and the repository
already contains the separating example.
$\mathrm{allEqual}(n,\alpha)$ is generated by the $\binom n 2$ (or
even $n-1$ consecutive) pairwise-equality constraints --- scopes of
size $2$, so $\mathrm{gw} = 1$ --- while its primal graph is complete
and its induced width is maximal ($n - 1$; prose observation in
\path{Ste.CouplingLowerBound}, not a Lean theorem).  More brutally:
graph $3$-colorability is generated by edge-scoped disequality
constraints, so it \v{C}ech-vanishes over the edge cover by
Proposition~\ref{prop:supportcover}(a) (this instantiation is prose,
not Lean); deciding whether $T \ne \emptyset$ is still NP-hard.  \emph{Vanishing of the gluing
obstruction over the natural cover is therefore compatible with full
computational hardness.}  All the hardness has been pushed into
computing the objects the proposition quantifies over --- the genuine
local sections $\Sec_T(U_j)$ --- and none remains in the gluing.
\end{observation}

We regard Observation~\ref{obs:arity} as the honest answer to ``is
low $\Hech^1$ over a good cover an algorithm?''  It is half of one:
\emph{gluing is the free half; sections are the expensive half}.  The
next section shows that the repository's algorithmic layer is
precisely a theory of when the expensive half is cheap.

\section{Reading the algorithmic layer as section-presheaf
computation}
\label{sec:algorithms}

\subsection{Bag tables are local sections}

The treewidth chain (\path{Ste.Treewidth} through
\path{Ste.JunctionTree}) is built on the \path{table} of a constraint
on a scope $\sigma$: the image of $T$ under restriction to $\sigma$.
Comparing definitions, \path{table}~$\sigma$~$T$ and
\path{localSections}~$T$~$\sigma$ are the \emph{same construction} ---
both are the image of $T$ under $f \mapsto (v : \sigma) \mapsto f\,v$;
the identification lemma is a one-line \texttt{rfl}-adjacent statement
that simply has not been written down.  So the machine-checked
junction-tree theorems are, without change, theorems about the section
presheaf on the bag cover:

\begin{itemize}
\item \emph{Faithfulness is the $\Hech^0$ statement}: a scoped
  constraint is exactly the preimage of its bag table
  (\path{HasSupport.eq_preimage_table}), and a global assignment is
  feasible for the instance iff it restricts into every bag table
  (\path{mem_joinConstraint_iff_restrict},
  \path{junctionTables_faithful}) --- membership in $T$ is decided by
  the family of its local-section sets over the bag cover.  This is
  Proposition~\ref{prop:supportcover}(b) in concrete clothing.
\item \emph{Cost}: a bag of at most $w+1$ variables over alphabets of
  size at most $k$ tables in at most $k^{w+1}$ rows
  (\path{table_encard_le_pow}); a width-$w$ elimination order
  materializes at most $n \cdot k^{w+1}$ rows in total
  (\path{bucketEliminate_total_space},
  \path{junctionTree_size_le}, linear-in-$n$ form
  \path{junctionTree_size_linear} at the minimal duplicate-free
  width \path{junctionWidth}).
\item \emph{Section computation}: projective bucket elimination
  computes exactly the projection of the joint solution set onto the
  remaining variables at each step
  (\path{joinConstraint_projectBucketEliminate} and neighbours in
  \path{Ste.BucketConsistency}), decides feasibility
  (\path{projectBucketEliminate_decides},
  \path{bucketEliminate_decides}), and extracts a solution
  backtrack-free
  (\path{exists_extension_of_mem_projectBucketEliminate}); the fused
  statement is \path{projectBucketEliminate_complete_solver} and, for
  adaptive consistency proper, \path{decreasingRun_adaptiveSolver}
  \citep{dechter1999bucket,dechter2003constraint,dechterpearl1987network}.
\end{itemize}

So the mechanized algorithmic layer already computes the genuine
section presheaf over an elimination cover, in $O(n \cdot k^{w+1})$
at induced width $w$ (\path{inducedTreewidth},
\path{achievesWidth_inducedTreewidth}; the two-sided bridge to
Robertson--Seymour treewidth \citep{robertson1986graphminorsII} is
\path{treewidth_primalGraph_le} and
\path{inducedTreewidth_le_treewidth_sub_one}; nonserial dynamic
programming ancestry in \citet{bertele1972nonserial}).

\subsection{Consistency algorithms compute pseudo-sections}

Local-consistency algorithms
\citep{mackworth1977consistency,freuder1982backtrack} compute
something weaker than $\Sec_T$: tables that are \emph{locally}
consistent but not necessarily globally extendable.  Call these
\emph{pseudo-sections}.  In Abramsky--Brandenburger vocabulary a
pseudo-section presheaf is exactly an empirical model over the cover
\citep{abramsky2011sheaf}; contextuality --- compatible local data
with no global explanation --- is the generic situation for
pseudo-sections, and its cohomological detection is the AMB program
\citep{abramsky2012cohomology,abramsky2015paradox,caru2017cohomology}.
The mechanized tractability results are precisely theorems that
\emph{pseudo-sections are genuine} under structural hypotheses:

\begin{itemize}
\item on a chain, arc consistency makes every domain value lie on a
  global solution (\path{arcConsistent_backtrackFree},
  \path{arcConsistent_chain_solvable}; the file's own docstring
  already reads this as gluing over the edge cover, whose nerve is an
  acyclic path);
\item on a rooted tree, directional arc consistency does the same
  (\path{treeArcConsistent_solvable}, \path{Ste.ConsistencyTree});
\item at induced width $w$, directional consistency of buckets does
  the same (\path{directionalConsistent_solvable},
  \path{Ste.AdaptiveConsistency}), and bucket elimination
  \emph{establishes} the condition while preserving solutions
  (\path{Ste.BucketConsistency}).
\end{itemize}

This gives a clean division of labour between the two theories:

\begin{center}
\begin{tabular}{@{}lll@{}}
\toprule
 & pseudo-sections ($k$-consistent tables) & genuine sections
 (\path{localSections}) \\
\midrule
cheap to compute & yes (local closure) & only at bounded width \\
per-section obstruction & can be nonzero (contextuality) & always
zero (\path{amb_extension_always}) \\
family gluing & fails generically & iff cover refines a generating
arity (Prop.~\ref{prop:supportcover}) \\
where hardness sits & pseudo $\ne$ genuine gap & computing
$\Sec_T$ \\
\bottomrule
\end{tabular}
\end{center}

\noindent
The conjectural summary: \emph{the $\Hech^1$-shaped content of STE is
the pseudo/genuine gap, and every tractability theorem in the
repository is a vanishing theorem for it over a structured cover.}
Making ``pseudo-section presheaf'' a first-class Lean object and
proving ``RIP cover + directional consistency $\Rightarrow$ pseudo
$=$ genuine'' as a single presheaf-level statement (subsuming the
chain, tree, and adaptive cases) is mechanization target~2 in
Section~\ref{sec:mechanization}.

\section{Quantitative invariants: obstruction counts, rectangle
covers, and what is genuinely incomparable}
\label{sec:quant}

\path{Ste.RepresentationBounds} settles --- negatively, and
machine-checked --- the tempting slogan ``obstruction size $=$ forced
representation blow-up.''  The rectangle-cover number $\rectrho(T)$
(\path{rectCoverNumber}): the least number of rectangles whose union
is exactly $T$, the analogue of the nondeterministic
communication-complexity cover number $C^1$
\citep{kushilevitz1997communication,aho1983notions}, with the
LP-extension-complexity lineage of
\citet{yannakakis1991expressing,fiorini2015exponential} explicitly
flagged as outlook.  Machine-checked:

\begin{itemize}
\item the fooling-set lower bound and the cardinality upper bound,
  $\mathrm{fool}(T) \le \rectrho(T) \le |T|$
  (\path{foolingNumber_le_rectCoverNumber},
  \path{IsFoolingSet.encard_le_of_hasRectCover},
  \path{rectCoverNumber_le_encard}); the mixing engine is
  \path{IsRectangle.mix_mem};
\item the bottom of the scale is cohomological: $\rectrho(T) = 1
  \iff$ the singleton-cover obstruction vanishes
  (\path{rectCoverNumber_eq_one_iff_cechVanishes},
  \path{one_lt_rectCoverNumber_iff});
\item on the coupling, $\rectrho(\mathrm{allEqual}(n,\alpha)) =
  |\alpha|$ exactly (\path{rectCoverNumber_allEqual},
  fooling set \path{isFoolingSet_allEqual},
  \path{foolingNumber_allEqual});
\item \emph{no uniform inequality holds in either direction} between
  the raw obstruction magnitude and $\rectrho$
  (\path{not_forall_cechObstruction_le_rectCoverNumber},
  \path{not_forall_rectCoverNumber_le_cechObstruction}): rectangles
  give $\obs = 0 < 1 = \rectrho$, and $\mathrm{allEqual}(3,
  \mathrm{Bool})$ gives $\rectrho = 2 < 6 = \obs$
  (\path{rectCoverNumber_lt_cechObstruction_allEqual}); at the
  minimal nonvanishing instance the two coincide
  (\path{diagonal_cechObstruction_eq_rectCoverNumber}) --- a
  coincidence of the smallest example, not a law.  The only equality
  that survives is the tautological box-gap reading
  (\path{cechObstruction_eq_box_sub_encard}).
\end{itemize}

Two representation grammars are in play and should be kept apart:
$\rectrho$ measures $\cup$-of-boxes (DNF-like) representations;
the junction-tree machinery measures $\cap$-of-cylinders
(CNF-over-bags) representations.  The machine-checked triple
$(\obs, \rectrho, \text{width})$ on $\mathrm{allEqual}$ shows all
three can diverge: obstruction exponential, $\rectrho$ constant,
width maximal.

\begin{conjecture}[$\rectrho$ vs.\ width incomparability, other
direction; not mechanized]
A width-$1$ chain of disequality constraints over an alphabet of size
$k \ge 3$ on $n$ variables has feasible set of size
$k(k-1)^{n-1}$ and, we conjecture, $\rectrho = 2^{\Omega(n)}$ (a
fooling-set construction along alternating assignments looks
plausible).  If so, $\rectrho$ and treewidth are incomparable in both
directions, and neither is a substitute for the other as a
``compressibility'' certificate: $\rectrho$ answers ``how many boxes,''
width answers ``how large a bag.''
\end{conjecture}

\emph{Practical reading for the STE practitioner}: fooling sets are
cheap-to-exhibit certificates that no small union-of-boxes
representation of a shrunken feasible family exists; stuck compatible
families over a proposed cover (Corollary~\ref{cor:scopebound}) are
cheap-to-exhibit certificates that no generating family respects
those bags.  Both certificate kinds are fully mechanized at the
definition level and computed exactly on the coupling family.

\section{The corpus axis, and the two-axis picture}
\label{sec:corpus}

The dynamic-frame gluing (\path{Ste.DynamicFrame.PresheafGluing})
deserves to be read side by side with
Proposition~\ref{prop:supportcover}, because it is the same theorem
shape on a different index: for corpora $D, E$, a hypothesis is
feasible on $D \cup E$ iff feasible on both pieces and the
\emph{straddling} constraints --- support in $D \cup E$ but in
neither piece --- are satisfied (\path{mem_feasible_union_iff}, with
the straddlers isolated as \path{crossObstruction}).  That is
exactly the two-context case of the support--cover analysis
\emph{with the obstruction term left explicit} instead of assumed
away by scope refinement: when no constraint straddles, gluing is
unconditional; when some do, the obstruction is a finite, per-candidate
checkable set.  Uniqueness of the glued world is a subsingleton
instance (\path{globalExtension_subsingleton}) --- the corpus axis'
compatibility notion (identical ambient hypothesis) makes it a
\emph{separated} presheaf, which is why its gluing theory closes at
$\Hech^0$ with no higher obstruction.

Conjecturally (not mechanized, and at this point frankly
speculative): the right home for both axes is a single bi-indexed
presheaf on $\text{(corpora)} \times \text{(variable contexts)}$,
contravariant in both, with the corpus direction separated and the
variable direction carrying the twisted cohomology of
Section~\ref{sec:coefficients}.  A Mayer--Vietoris-style organization
of iterated corpus splits (the standard machinery, e.g.\
\citealp{bott1982forms}) would give divide-and-conquer over
provenance with correction terms that are precisely the
\path{crossObstruction} sets; since the corpus presheaf is separated,
we expect the ``spectral sequence'' to degenerate immediately, which
is a polite way of saying the corpus axis needs no cohomology beyond
what is already mechanized --- but a mechanized statement of the
$k$-piece inclusion--exclusion form of \path{mem_feasible_union_iff}
would still be a useful lemma for shard-parallel STE, and the
variable-side analogue with an explicit obstruction term (splice two
partial solutions, collect the straddling constraints) is
mechanization target~3 below.

One more cross-connection worth recording: on the corpus axis the
practical solvers (\path{Ste.DynamicFrame.FiniteSolver}, the
TIDES/Federalist corpus instantiations in \path{Ste.TidesCorpus},
\path{Ste.FiniteInstance}, \path{Ste.TidesComputable}) work
extensionally and hit the $2^N$ wall by design; the variable-side
machinery of this memo is exactly what a second-generation solver
would use to stay under it when the claim-interaction hypergraph has
bounded arity and width.  Carlson's cipher instantiation
\citep{carlson2012} sits at the opposite corner: there the hypothesis
space (keys) is unstructured but the property sets are highly
structured, and the productive reductions are class-level embeddings
(\path{permutation_reduces_to_substitution}) rather than
variable-local gluing.

\section{Three speedup routes, ranked, and the mechanization queue}
\label{sec:mechanization}

\subsection{The three routes}

\textbf{Route 1 (essentially proven; assembly work remains): bag-local
query answering.}  At induced width $w$ over alphabet size $k$:
compute the genuine section presheaf over the elimination cover by
projective bucket elimination ($O(n \cdot k^{w+1})$ total table
space, machine-checked as cited in Section~\ref{sec:algorithms});
answer feasibility (\path{projectBucketEliminate_decides}),
extract witnesses backtrack-free, and answer must/may/cannot pair
queries whose variables co-occur in a bag directly from that bag's
section table --- soundness is faithfulness
(\path{mem_joinConstraint_iff_restrict}) plus, for the gluing step,
Proposition~\ref{prop:supportcover}(a), now machine-checked
(\path{cechVanishesCover_of_refiningGenerators}).  For pairs not
co-occurring in a bag, propagate through separators; this is standard
junction-tree practice \citep{dechter2003constraint} but the
composite correctness statement (``the bag-local answer equals the
global must/may answer'') is \emph{not yet a Lean theorem} --- it is,
however, within reach of the existing lemmas, and it is the statement
that turns the $2^N$-wall companion note's Section on structural
compression into a query engine.  Expected speedup: from $2^N$
extensional enumeration to $O(n \cdot k^{w+1})$ preprocessing plus
$O(k^{w+1})$ per query, on bounded-width instances.

\textbf{Route 2 (conjectural): the linear-gluing relaxation as a
cheap refutation engine.}  \path{LinearlyGlues} is a linear-algebra
condition; its failure refutes set-level gluing
(contrapositive of \path{GluesCover.linearlyGlues}, machine-checked),
and it is strong enough to catch the coupling over every nontrivial
ring (\path{mixedFamily_not_linearlyGlues}).  As stated, the linear
system ranges over global solutions and is useless
computationally; the conjecture is that over a bag cover one can
present a \emph{relaxed} system over per-bag pseudo-sections with
consistency constraints on separators --- a signed / affine analogue
of the $k$-consistency hierarchy, cousin to Sherali--Adams-style CSP
relaxations and to the contextual-fraction linear programs of
\citet{abramsky2017contextualfraction} --- whose infeasibility is a
certificate of STE infeasibility computable in time polynomial in the
table sizes.  What vanishing of the twisted $\Hech^1$ over structured
covers would add is exactness guarantees for such relaxations; the
open \path{TwistedH1Trivial} question for disjoint covers is the
first stepping stone.  Nothing here is mechanized; even the right
statement is open, and cohomological-vs-contextual comparisons in the
literature \citep{caru2017cohomology,okay2017topological,
aasnaess2022comparing} caution that $\Hech^1$-style invariants can
have false negatives --- indeed the mechanized
\path{amb_extension_always} is precisely such a false-negative
theorem.

\textbf{Route 3 (proven ingredients; new composite): obstruction
certificates before you pay for elimination.}  Before committing to
an elimination order, test candidate covers: a single stuck
compatible family over $U$ (checkable on small witnesses; the
repository computes them exactly for the coupling family) certifies
via Corollary~\ref{cor:scopebound} that no generating family refines
$U$, so no $U$-bag-local pipeline is sound; a fooling set certifies
no small $\cup$-of-boxes representation
(\path{IsFoolingSet.encard_le_of_hasRectCover}).  Cheap NO-certificates
steer the expensive machinery; this is the cohomology-as-diagnostics
reading, and it is available today.

\subsection{The mechanization queue}

The original draft of this memo put the support--cover proposition at
the head of this queue; \textbf{that item is now done} ---
Proposition~\ref{prop:supportcover} is mechanized in
\path{Ste.SupportCover}
(\path{cechVanishesCover_of_refiningGenerators},
\path{cechVanishesCover_iff_eq_iInter_sectionCylinder},
\path{cechVanishesCover_iff_exists_scopedGeneration}).  The updated
queue, in descending expected value per line of Lean:

\begin{enumerate}
\item \textbf{Harvest the proposition's corollaries.}  The one-line
  \path{table} $=$ \path{localSections} identification lemma
  (Section~\ref{sec:algorithms}); a Lean definition of the gluing
  width (Definition~\ref{def:gluingwidth}; no such definition exists
  yet) with the
  arity characterization read off the mechanized iff; and
  Corollary~\ref{cor:scopebound} as its contrapositive.  All three
  are short, and together they turn the new theorem into the
  certificate interface of Route~3.
\item \textbf{The pseudo-section presheaf} and the vanishing theorem
  ``RIP cover $+$ directional consistency $\Rightarrow$
  pseudo-sections are genuine,'' unifying
  \path{arcConsistent_backtrackFree},
  \path{treeArcConsistent_solvable}, and
  \path{directionalConsistent_solvable} as one presheaf statement.
  This is where the genuinely $\Hech^1$-shaped content of STE lives
  (Section~\ref{sec:algorithms}), and it would make the
  Abramsky--Brandenburger comparison exact rather than analogical.
\item \textbf{Variable-side cross-obstruction gluing}: the splice of
  two partial solutions over $W_1 \cup W_2$ with an explicit
  straddler-obstruction term, mirroring
  \path{mem_feasible_union_iff} on the variable axis; then its
  $k$-piece iteration along an RIP ordering.  This is the
  Mayer--Vietoris lemma of Route~1's separator propagation.
\item \textbf{The bag-local query theorem} of Route~1 (must/may/cannot
  correctness from junction tables), and the quantitative
  lower-bound side the junction-tree file lists as outlook (a
  width-$w$-vs-$2^{\Omega(n)}$ separation, for which
  \path{cechObstruction_allEqual_ge} is the margin-cover prototype).
\item \textbf{Twisted $\Hech^1$ for genuine covers}: resolve the
  quotient-type diamond noted in \path{Ste.TwistedCech}, settle
  \path{TwistedH1Trivial} for disjoint covers, and only then attempt
  the connecting map $\gamma$ of \citet{abramsky2012cohomology} ---
  with the caveat that \path{amb_extension_always} already proves the
  per-section $\gamma$ is identically blind in STE, so the payoff is
  the \emph{family-level} cokernel theory, not $\gamma$ itself.
\end{enumerate}

\section{Summary}

What is machine-checked: the constraint side is a sheaf
(\path{feasibilitySet_eq_iInter_cover}); the variable-side section
presheaf is a functor whose singleton-cover gluing characterizes
rectangularity exactly (\path{cechVanishes_iff_rectangular}), whose
arbitrary-cover gluing holds for rectangles with unique glues
(\path{rectangular_glueCover_existsUnique}) and fails at couplings
(\path{diagonal_not_cechVanishesCover}) with exactly computed,
exponentially growing obstruction counts
(\path{cechObstruction_allEqual}, \path{cechObstruction_allEqual_ge});
constant-coefficient \v{C}ech cohomology of the full nerve is
provably blind (\path{cechH1_subsingleton}) while the twisted
degree-$0$ cokernel class fires on the coupling over every nontrivial
ring (\path{mixedFamily_not_linearlyGlues}) even though the AMB
per-section $\Hech^1$ obstruction provably cannot
(\path{amb_extension_always}); the rectangle-cover number is
sandwiched, characterized at $1$ cohomologically, computed on the
coupling, and provably incomparable to the obstruction magnitude
(\path{not_forall_cechObstruction_le_rectCoverNumber} and its
mirror); and the entire bounded-width algorithmic stack --- tables,
bucket elimination, junction trees, adaptive consistency --- is
mechanized with faithfulness and $O(n \cdot k^{w+1})$ size and a
fused complete-solver statement
(\path{projectBucketEliminate_complete_solver}); and --- new since the
first draft of this memo --- the support--cover gluing theorem
itself: gluing of genuine sections is governed entirely by scope
refinement (Proposition~\ref{prop:supportcover}, mechanized in
\path{Ste.SupportCover}:
\path{cechVanishesCover_of_refiningGenerators},
\path{cechVanishesCover_iff_eq_iInter_sectionCylinder},
\path{cechVanishesCover_iff_exists_scopedGeneration}).

What we still conjecture: the downstream reading of that theorem ---
the scope lower bound, the gluing-width invariant, and the
arity-vs-treewidth separation
(Corollary~\ref{cor:scopebound},
Definition~\ref{def:gluingwidth},
Observation~\ref{obs:arity}) --- remains sketch-level; that the
cohomological content of STE hardness is the
pseudo-section/genuine-section gap, with every mechanized
tractability theorem a vanishing theorem for it; that a bag-local
linear relaxation with exactness controlled by twisted cohomology
exists; and that $\rectrho$ and width are incomparable in both
directions.  The single highest-value next mechanization is no longer
the proposition --- that is done --- but its harvest: the
\path{table} $=$ \path{localSections} identification and a Lean-level
gluing-width definition with Corollary~\ref{cor:scopebound}, which
together make the new theorem load-bearing for the only speedup route
fully within reach of the current development.  Pleasingly for a memo
about cohomology, the theorem's content is that in set-theoretic
estimation the cohomology, properly aimed, is exactly the bookkeeping
of who is allowed to talk to whom --- and that bookkeeping is now
machine-checked.

\bibliographystyle{plainnat}
\bibliography{../../refs}

\end{document}
